\documentclass[11pt]{article}

\usepackage[a4paper,margin=2.5cm]{geometry}
\usepackage{amsmath,amssymb,amsthm}
\usepackage[colorlinks=true,linkcolor=blue,citecolor=blue,urlcolor=blue]{hyperref}
\usepackage{booktabs}
\usepackage{enumitem}
\usepackage{natbib}

\newtheorem{theorem}{Theorem}[section]
\newtheorem{lemma}{Lemma}[section]
\newtheorem{proposition}{Proposition}[section]
\newtheorem{corollary}{Corollary}[section]
\newtheorem{definition}{Definition}[section]
\newtheorem{remark}{Remark}[section]
\newtheorem{hypothesis}{Hypothesis}[section]

\theoremstyle{remark}
\newtheorem*{importedTheorem}{Imported theorem (citation only)}

\newcommand{\twoI}{2I}
\newcommand{\Vsixhundred}{V_{600}}
\newcommand{\Dicfive}{\mathrm{Dic}_5}
\newcommand{\Vtwentyfour}{V_{24}}
\newcommand{\Sigmaspine}{\Sigma_{\Vsixhundred}}
\newcommand{\CQ}{C_{\mathbb{Q}}}
\newcommand{\Ttau}{T_{\tau}}
\newcommand{\sigmaG}{\sigma}
\newcommand{\tausigma}{\tau_\sigma}
\newcommand{\AK}{\mathcal{A}_K}
\newcommand{\PK}[1]{P_{#1}}
\newcommand{\Hzero}{H_0}

\title{A unified $\Vsixhundred$ scaffolding for finite-group derivations of\\
       $1/4$ incidence, monochromatic $\sigma$-pair energy,\\
       Galois-coset duality, and operator-trace ratios}
\author{%
  Lee Smart \\[3pt]
  \textit{Institute of Vibrational Field Dynamics} \\[2pt]
  \texttt{contact@vibrationalfielddynamics.org}}
\date{May 2026}

\begin{document}

\maketitle

\noindent\textbf{Status:} Preprint (not peer-reviewed).
This is a SYNTHESIS paper. It cites four prior papers in the
$\Vsixhundred$ programme and adds NO new theorem-grade physical
claim. Its sole contribution is the explicit identification of a
shared structural tuple $\Sigmaspine$ that all four imported
theorems factor through.

\begin{abstract}
We synthesise four prior results on the binary icosahedral group
$\Vsixhundred = \twoI$:
the Bekenstein--Hawking $1/4$ coefficient as $\Dicfive$-coset
incidence on $\Vsixhundred$~\citep{V600Paper1};
the discrete Hawking quantum $E_q = 5/2$ from $\sigmaG$-pair
excitations on $\Vsixhundred$~\citep{V600Paper2};
the canonical involution $\tausigma : \Vsixhundred \to \Vsixhundred$
with explicit phase-$0$ lift and $\mathbb{Z}_2^5 = 32$ phase
ambiguity~\citep{V600Paper3};
and the exact rational operator-trace ratios
$\mathrm{tr}(I_{\CQ} + \PK{72})/12 = 13/12$ and
$\mathrm{tr}(I_{\CQ} - \PK{0})/12 = 11/12$ on the
$12$-dimensional rational $\Ttau$-cycle space, with the
cosmological coupling rule stated as an explicit Layer-3
hypothesis~\citep{V600Paper4}.
We record the load-bearing structural tuple
\[
  \Sigmaspine \;=\; \bigl(\,\twoI,\; \Ttau\text{-cycles},\;
  \text{K-multiset } \{72{:}1,\,0{:}1,\,52{:}5,\,20{:}5\},\;
  \Dicfive,\; \Vtwentyfour,\;
  \text{right cosets } \Dicfive g,\;
  \sigmaG,\; \tausigma,\; \AK\,\bigr),
\]
and show that all four theorem-grade imports, with Paper~4
restricted to its Layer~1 trace identities and admissibility
theorem, factor through it.
The synthesis adds NO new theorem-grade physical claim. We do NOT
claim a unified physical theory; we report a unified mathematical
scaffolding. CMB-bulk projection, cosmic dipole ladder audits, and
pre-registered prediction manifests are explicitly named future
builds with unmet prerequisites; they are out of scope here.
\end{abstract}

\section{Introduction}\label{sec:intro}

The $\Vsixhundred$ programme is five papers. Papers~1--4 each prove
a single narrow finite-group identity:
\begin{itemize}[leftmargin=2em]
  \item Paper~1~\citep{V600Paper1}: per-coset incidence count
  $|C \cap \Vtwentyfour| = 4$ for $C \in \{g\Dicfive, \Dicfive g\}$,
  hence $S/A = 4/16 = 1/4$.
  \item Paper~2~\citep{V600Paper2}: $|v - \sigmaG(v)|_E^2 = 5/2$ for
  every $\sigmaG$-mobile vertex of $\Vsixhundred$, hence a single-line
  $\sigmaG$-pair excitation spectrum.
  \item Paper~3~\citep{V600Paper3}: explicit $\tausigma$ involution
  fixing $\Dicfive$ pointwise, swapping $K = 52$ and $K = 20$ cycle
  classes within each non-bulk right coset, and commuting with
  $\Ttau$ and the antipodal map.
  \item Paper~4~\citep{V600Paper4}: $\mathrm{tr}(I_{\CQ} + \PK{72}) = 13$
  and $\mathrm{tr}(I_{\CQ} - \PK{0}) = 11$ on the $12$-dimensional
  cycle space, plus a K-saturated admissibility theorem and an
  explicit Layer-3 cosmological coupling hypothesis.
\end{itemize}

The present paper is the synthesis. It does NOT prove anything new
about physical reality. It identifies the cross-cutting finite
structure that makes all four narrow theorems work, and records that
identification as one structural-spine table plus four citation-only
imported-theorem boxes. Anything that would require new theorems
about cosmological projection, observed CMB, or large-scale-survey
predictions is named out-of-scope (Section~\ref{sec:scope}).

\section{Structural spine}\label{sec:spine}

We collect the load-bearing finite data into one tuple.

\begin{definition}[Structural spine]\label{def:spine}
The \emph{structural spine} of the $\Vsixhundred$ programme is the
finite tuple
\[
  \Sigmaspine \;=\; \bigl(\,
    \twoI,\;
    \Ttau,\;
    \mathrm{K\text{-multiset}},\;
    \Dicfive,\;
    \Vtwentyfour,\;
    \mathcal{R},\;
    \sigmaG,\;
    \tausigma,\;
    \AK\,
  \bigr),
\]
whose components are the load-bearing facts of
Table~\ref{tab:spine}. All entries are imported from the cited
foundation papers; no entry is introduced or re-proved here.
\end{definition}

\begin{table}[h]
\centering
\caption{Components of the structural spine $\Sigmaspine$. Every
entry is imported from the cited paper or from the shared library
\texttt{vfd\_v600}; this paper does not re-prove any of them.}
\label{tab:spine}
\renewcommand{\arraystretch}{1.15}
\begin{tabular}{p{3cm}p{6cm}p{3cm}p{2cm}}
\toprule
Component & Description & Source & Status \\
\midrule
$\twoI = \Vsixhundred$ & Binary icosahedral group; $120$ unit
icosian quaternions; vertices of the regular $600$-cell. &
\citet{Coxeter1973,ConwaySloane1988}, Paper~1 setup & exact \\
$\Ttau$-cycles & $T_\tau : g \mapsto \tau g$ partitions
$\Vsixhundred$ into $12$ cycles each of length $10$, for any
$\tau \in \twoI$ of order $10$. & Paper~1 & exact \\
K-multiset & The $12$ cycles carry K-class assignments
$\{72{:}1,\,0{:}1,\,52{:}5,\,20{:}5\}$. & Paper~1 & exact \\
$\Dicfive$ & Binary dihedral subgroup of order $20$; equals the
union of the $K=72$ and $K=0$ cycles; self-normalising and non-normal
in $\twoI$. & Paper~1 & exact \\
$\Vtwentyfour$ & $\sigmaG$-fixed sublattice
$\mathrm{Fix}_\sigma(\Vsixhundred)$; $24$ vertices; the
$24$-cell ($8$ axis $+$ $16$ half-type). & Paper~1 & exact \\
$\sigmaG$ & The $K/\mathbb{Q}$ Galois twist
$\sqrt{5} \mapsto -\sqrt{5}$ on the icosian quaternion algebra
$\mathbb{H}_K$, restricted to act on $\Vsixhundred$. Defines
$\Vtwentyfour = \mathrm{Fix}_\sigma$ and the $\sigmaG$-pair energy
$|v - \sigmaG(v)|_E^2$. & Paper~1, Paper~2 & exact \\
$\mathcal{R}$ & The five non-bulk right cosets $\Dicfive g$, each
the disjoint union of one whole $K{=}52$ cycle and one whole
$K{=}20$ cycle. Left cosets $g\Dicfive$ are NOT whole-cycle carriers.
& Paper~1, Paper~3 & exact \\
$\tausigma$ & Canonical involution
$\tausigma : \Vsixhundred \to \Vsixhundred$ with phase-$0$ lift
explicit on $120$ vertices; $\mathbb{Z}_2^5 = 32$ phase ambiguity;
fixes $\Dicfive$ pointwise; swaps $K=52 \leftrightarrow K=20$ within
each non-bulk right coset; commutes with $\Ttau$ and the antipodal
map. & Paper~3 & exact \\
$\AK$ & K-saturated diagonal projector algebra
$\langle I_{\CQ}, \PK{72}, \PK{0}, \PK{52}, \PK{20}\rangle_\mathbb{Q}$
on the $12$-dimensional rational cycle space $\CQ \cong \mathbb{Q}^{12}$.
& Paper~4 & exact \\
\bottomrule
\end{tabular}
\end{table}

\subsection{Correctness lemma: $\Vtwentyfour \neq \Dicfive$}

A subtle conflation in earlier exploratory drafts must be flagged.

\begin{lemma}[$\Vtwentyfour$ / $\Dicfive$ distinction; Paper 1]\label{lem:V24Dic5}
$\mathrm{Fix}_{\sigmaG}(\Vsixhundred) = \Vtwentyfour$ has $24$
vertices ($8$ axis $+$ $16$ half-type). $\Dicfive$ is the bulk
subgroup of order $20$ formed by the union of the $K=72$ and $K=0$
cycles. These are distinct subsets of $\Vsixhundred$;
$|\Vtwentyfour| = 24$, $|\Dicfive| = 20$, and
$|\Dicfive \cap \Vtwentyfour| = 4$. More generally,
$|C \cap \Vtwentyfour| = 4$ for every left or right coset
$C \in \{g\Dicfive, \Dicfive g\}$ --- this is the
Bekenstein incidence theorem (Paper~1, not a new claim of
this paper).
\end{lemma}

\begin{proof}
Imported from Paper~1~\citep{V600Paper1}. The $4$/$16$ ratio
underlies the Bekenstein $1/4$.
\end{proof}

\subsection{Phase-independence}

\begin{lemma}[Cross-cutting phase-independence; Papers 3, 4]\label{lem:phaseind}
The theorem-grade identities and properties imported below
(Sections~\ref{sec:p1}--\ref{sec:p4}) are independent of which
canonical $\tausigma$ lift (out of the $\mathbb{Z}_2^5 = 32$ canonical
lifts) is chosen. Specifically: the K-class projector traces, the
$|gH \cap \Vtwentyfour|$ count, the $\sigmaG$-pair-energy spectrum,
and the involution $(P1)$--$(P4)$ properties are all functions of
$\Sigmaspine$ that do not depend on the phase choice. (The explicit
phase-$0$ lift exhibited in Paper~3 is, by construction, phase-specific
as a $120$-vertex map; the structural properties it satisfies are
phase-independent.)
\end{lemma}

\begin{proof}
Phase-independence of $(P1)$--$(P4)$ is in
Paper~3~\citep{V600Paper3} (all $32$ lifts satisfy the same
properties). Phase-independence of the operator-trace ratios is
verified explicitly in Paper~4~\citep{V600Paper4} (all $32$ lifts
yield cycle-level permutations that pointwise fix $\hat H$ and
$\hat C$). The $1/4$ incidence count and the $\sigmaG$-pair
spectrum are properties of $\Vsixhundred$ and $\sigmaG$ alone and do
not reference $\tausigma$.
\end{proof}

\section{Imported theorem 1: Bekenstein \texorpdfstring{$1/4$}{1/4}
incidence (Paper 1)}\label{sec:p1}

\begin{importedTheorem}[\citealp{V600Paper1}]
For every coset $C \in \{g\Dicfive,\, \Dicfive g\}$ of the bulk
subgroup $\Dicfive \subset \twoI$, the intersection
$|C \cap \Vtwentyfour| = 4$, equivalently $4$
$\sigmaG$-fixed vertices and $16$ $\sigmaG$-mobile vertices per
coset (left or right). The structural ratio $S/A = 4/16 = 1/4$
matches the Bekenstein--Hawking entropy coefficient.
\end{importedTheorem}

\paragraph{Role in synthesis.}
Imported theorem~1 uses $(\Vsixhundred, \Dicfive, \Vtwentyfour)$
from $\Sigmaspine$. The result is purely arithmetic and does not
require $\tausigma$ or $\AK$. Paper~1 is the load-bearing setup
paper for the rest of the programme; its right-coset/left-coset
distinction (Lemma in Paper~1) is the load-bearing fact for
Papers~3 and 4.

\section{Imported theorem 2: Hawking quantum $E_q = 5/2$ (Paper 2)}\label{sec:p2}

\begin{importedTheorem}[\citealp{V600Paper2}]
Every $\sigmaG$-mobile vertex $v \in \Vsixhundred \setminus \Vtwentyfour$
satisfies $|v - \sigmaG(v)|_E^2 = 5/2$ in the Euclidean norm. The
$96$ $\sigmaG$-mobile vertices yield a monochromatic $\sigmaG$-pair
excitation spectrum $24 \cdot \delta_0 + 96 \cdot \delta_{5/2}$.
A finite first-law identity holds per coset (left or right; the bulk
and non-bulk cases share the same per-coset arithmetic):
$\mathcal{T}_{\mathrm{casc}} \cdot S = E/4$ with $A = 16$, $S = 4$,
$E = 40$, $\mathcal{T}_{\mathrm{casc}} = 5/2$, where
$\mathcal{T}_{\mathrm{casc}}$ is the finite cascade temperature
parameter of Paper~2 (NOT calibrated to a physical Hawking
temperature; see Paper~2's explicit caveat).
\end{importedTheorem}

\paragraph{Role in synthesis.}
Imported theorem~2 uses $(\Vsixhundred, \Vtwentyfour, \sigmaG,
\Dicfive)$ from $\Sigmaspine$. It is independent of $\Ttau$,
$\tausigma$, and $\AK$. The first-law identity reuses the $4$/$16$
incidence from Paper~1.

\section{Imported theorem 3: $\tausigma$ involution (Paper 3)}\label{sec:p3}

\begin{importedTheorem}[\citealp{V600Paper3}]
There exists an explicit involution $\tausigma : \Vsixhundred \to
\Vsixhundred$ such that
$(P1)$ $\tausigma \circ \tausigma = \mathrm{id}$;
$(P2)$ $\tausigma$ fixes $\Dicfive$ pointwise;
$(P3)$ $\tausigma$ swaps the $K=52$ and $K=20$ cycle classes within
each non-bulk right coset of $\Dicfive$;
$(P4)$ $\tausigma$ commutes with $\Ttau$ (i.e.,
$\tausigma \circ \Ttau = \Ttau \circ \tausigma$);
and additionally $\tausigma$ commutes with the antipodal map
$v \mapsto -v$. There is a $\mathbb{Z}_2^5 = 32$ phase ambiguity in
the canonical lift; the phase-$(0,0,0,0,0)$ representative is
explicit on $120$ vertices.
\end{importedTheorem}

\paragraph{Role in synthesis.}
Imported theorem~3 uses $(\Vsixhundred, \Ttau, \mathrm{K\text{-multiset}},
\Dicfive, \mathcal{R}, \sigmaG)$ from $\Sigmaspine$ and produces
$\tausigma$. The construction is via right cosets $\mathcal{R}$ and
the icosian trace metric; the phase-independence
Lemma~\ref{lem:phaseind} is the load-bearing fact downstream.

\section{Imported theorem 4 (Paper 4): three layers}\label{sec:p4}

Paper~4 maintains a strict three-layer credibility split. We import
\emph{only Layer~1} as a theorem (IT4); Layer~2 and Layer~3 are
brought across separately as numerical observation and explicit
hypothesis, neither of which is promoted by this synthesis.

\begin{importedTheorem}[\citealp{V600Paper4}, Layer 1 only]
On the $12$-dimensional rational cycle space
$\CQ \cong \mathbb{Q}^{12}$:
\begin{enumerate}[leftmargin=2em]
\item \emph{Trace identities.}
$\mathrm{tr}(I_{\CQ} + \PK{72}) = 13$ and
$\mathrm{tr}(I_{\CQ} - \PK{0}) = 11$, hence ratios
$13/12$ and $11/12$ as exact rationals.
\item \emph{K-saturated admissibility.} Within $\AK$, the only
rank-one projector corrections to $I_{\CQ}$ are
$\pm \PK{72}$ and $\pm \PK{0}$.
\end{enumerate}
\end{importedTheorem}

\paragraph{Layer 2 (observation).}
Paper~4 reports
$\frac{13}{12} \cdot \Hzero^{\mathrm{Planck}} \approx 72.97$~km/s/Mpc
versus SH0ES $73.04 \pm 1.04$ ($\approx 0.06\sigma$ residual), and
$\frac{11}{12} \cdot S_8^{\mathrm{Planck}} \approx 0.763$ versus
KiDS-1000 multi-probe $0.766^{+0.020}_{-0.014}$
($\approx 0.24\sigma$ directional residual using the lower KiDS
error). These are observational, not derivations.

\paragraph{Layer 3 (hypothesis).}
Paper~4 states the cycle-mode coupling hypothesis $\Gamma$
explicitly~(\citealp{V600Paper4}, H1). The synthesis carries this
forward as the same hypothesis; we do NOT promote it to a theorem
here.

\paragraph{Role in synthesis.}
Imported theorem~4 uses $(\Vsixhundred, \Ttau, \mathrm{K\text{-multiset}},
\AK)$ from $\Sigmaspine$. Phase-independence
(Lemma~\ref{lem:phaseind}) ensures the trace ratios do not depend on
the choice of canonical $\tausigma$ lift.

\section{Synthesis}\label{sec:synthesis}

\begin{proposition}[Common factorisation through $\Sigmaspine$]\label{prop:synthesis}
Each of imported theorems~1, 2, 3, and IT4-Layer~1 is a statement
about $\Sigmaspine$. Specifically:
\begin{itemize}[leftmargin=2em]
\item Theorem~1 (Bekenstein $1/4$): a function of
$(\Vsixhundred, \Dicfive, \Vtwentyfour)$.
\item Theorem~2 (Hawking $E_q = 5/2$): a function of
$(\Vsixhundred, \Vtwentyfour, \sigmaG, \Dicfive)$ with the inherited
Euclidean icosian norm on the underlying $\mathbb{H}_K$.
\item Theorem~3 ($\tausigma$ \emph{structural-property package}
$(P1)$--$(P4)$ + antipodal compatibility): a function of
$(\Vsixhundred, \Ttau, \mathrm{K\text{-multiset}}, \Dicfive,
\mathcal{R}, \sigmaG, \tausigma)$. Note: the explicit phase-$0$
vertex map of Paper~3 also requires the icosian trace metric and
the deterministic base-point/ordering conventions of Paper~3's
\texttt{vfd\_v600.icosian.build\_vertices()} as auxiliary data;
those conventions are not part of $\Sigmaspine$ but are not needed
for the structural-property package itself.
\item IT4-Layer~1 (trace ratios + K-saturated admissibility): a
function of
$(\Vsixhundred, \Ttau, \mathrm{K\text{-multiset}}, \AK)$.
\end{itemize}
None of the four imported Layer~1 theorems uses any datum outside
$\Sigmaspine$. The factorisation does NOT extend to IT4 Layer~2
(observational benchmarks $\Hzero$, $S_8$ from external cosmology
data) or IT4 Layer~3 (the coupling hypothesis $\Gamma$ relating the
trace ratios to those observables); both Layer~2 and Layer~3
explicitly invoke data and physical hypotheses outside $\Sigmaspine$.
\end{proposition}

\begin{proof}
By inspection of each cited paper. The full derivations are in
\citet{V600Paper1,V600Paper2,V600Paper3,V600Paper4}; this paper
re-proves nothing.
\end{proof}

\begin{remark}[What the synthesis does and does not say]\label{rem:synth}
Proposition~\ref{prop:synthesis} is a structural observation: four
finite-group results share a common load-bearing structure. It is
NOT a unification theorem in physics; it does not derive black-hole
thermodynamics, Hawking radiation, or cosmological tensions. It
records that all four
identities depend on the same underlying tuple. Whether this
shared dependence has physical content is a separate question that
this paper does not address.
\end{remark}

\section{Honest scope and named future builds}\label{sec:scope}

The credibility-bar discipline of Papers~1--4 is preserved. We
explicitly name three classes of work that would extend the
synthesis but are NOT claims of this paper:

\paragraph{(F1) CMB-bulk projection bridge.}
The closure-cosmogenesis exploratory note in this repository
(\texttt{docs/closure-cosmogenesis.md}) contains conditional residual
and coarse-graining ingredients --- a conditional bulk-invariance
property (residual support $\rho^B_t = 0$ is conditional), an
unspecified coarse-graining $\pi_C$ with open weight choices, and no
temperature calibration --- that a future bridge map
$\Theta : (F^B_t, \pi_C, \text{calibration}) \to T : S^2 \to \mathbb{R}_+$
from a candidate bulk-supported sector of states (with
$\sigmaG$-invariance to be proved by the bridge, not assumed) to an
observed CMB temperature field would have to complete. The source
note does NOT contain such a map. The bulk-supported sector here is
NOT identified with $\Vtwentyfour$: the bulk subgroup $\Dicfive$ has
$20$ elements and $\Vtwentyfour$ has $24$; the two are distinct
(Lemma~\ref{lem:V24Dic5}), with $|\Dicfive \cap \Vtwentyfour| = 4$.
We do NOT use any of this in the synthesis; ``uniform
$T_{\mathrm{CMB}}$ follows from bulk invariance'' is NOT a theorem
of this paper, and any source language identifying the $20$-vertex
bulk with the $24$-vertex $\sigmaG$-fixed sublattice is a known
conflation that this synthesis explicitly disowns.

\paragraph{(F2) Cosmic dipole / $1+k/2$ ladder audit.}
Repository scripts catalogue a ladder of cascade-resolution
predictions and report agreement with several survey dipole-axis
measurements. Those scripts are exploratory and partly fitted; the
$k$-rule in
\texttt{papers/cosmological-folding-rate/dynamics/preregister\_k.py}
is a working hypothesis, not a frozen pre-registration. We do NOT
claim any dipole result in the synthesis; a separate empirical
audit with frozen survey inclusion, uncertainties, $k$-assignment,
and null model is needed before any such claim is defensible.

\paragraph{(F3) Pre-registered prediction manifest.}
A locked, timestamped prediction manifest with formula, value,
uncertainty window, and source commit per observable would convert
Hypothesis~$\Gamma$ (Paper~4) into a falsifiable programme. The
present scripts (\texttt{preregister\_k.py},
\texttt{cmb\_mapping.py}, the legacy
\texttt{UNIFIED\_RESULT.md}) are not in that form. The synthesis
does NOT claim a prediction catalogue; we cite the existing
exploratory artifacts only as programme context.

\paragraph{What the synthesis claims, in one sentence.}
There is a single finite tuple $\Sigmaspine$ such that the four
imported theorems factor through it. That is the entire claim.

\section{Conclusion}\label{sec:concl}

The four foundation papers of the $\Vsixhundred$ programme each
prove a narrow finite-group identity. This synthesis records that
they all factor through one structural tuple $\Sigmaspine$. No
new theorem-grade physical claim is added; CMB, dipole, and
prediction-manifest work is named out-of-scope with explicit unmet
prerequisites.

\appendix

\section{Verification certificate}\label{app:verify}

The companion script \texttt{verify.py} computes $\Sigmaspine$ from
\texttt{vfd\_v600.group.build\_state()} (using only the Python
standard library and the shared $\mathtt{vfd\_v600}$ package), and
asserts the listed headline cross-checks:
\begin{enumerate}[leftmargin=2em]
\item $|\Vsixhundred| = 120$, $12$ $\Ttau$-cycles each of length
$10$, K-multiset $\{72{:}1,\,0{:}1,\,52{:}5,\,20{:}5\}$.
\item $|\Dicfive| = 20$ and $|\Vtwentyfour| = 24$, so
$\Dicfive \neq \Vtwentyfour$ (Lemma~\ref{lem:V24Dic5}).
\item Bekenstein cross-check:
$|C \cap \Vtwentyfour| = 4$ for every left or right coset
$C \in \{g\Dicfive, \Dicfive g\}$; ratio $4/16 = 1/4$.
\item Hawking cross-check:
$|v - \sigmaG(v)|_E^2 = 5/2$ for every $\sigmaG$-mobile vertex;
spectrum is $24\cdot \delta_0 + 96\cdot\delta_{5/2}$.
\item $\tausigma$ cross-check (canonical phase-$0$ lift):
$(P1)$--$(P4)$ properties + antipodal compatibility hold.
\item Trace-ratio cross-check:
$\mathrm{tr}(I_{\CQ} + \PK{72}) = 13$ and
$\mathrm{tr}(I_{\CQ} - \PK{0}) = 11$ on $\CQ \cong \mathbb{Q}^{12}$.
\item Phase-independence: for each of the $2^5 = 32$ canonical
$\tausigma$ phase lifts, the induced cycle-level action pointwise
fixes both $I_{\CQ} + \PK{72}$ and $I_{\CQ} - \PK{0}$.
\end{enumerate}
Reproducibility command:
\begin{verbatim}
PYTHONPATH=papers/v600-programme/lib \
    python3 papers/v600-unified/verify.py
\end{verbatim}

\bibliographystyle{plainnat}
\bibliography{references}

\end{document}
